\documentclass[12pt]{article}
\usepackage{amsmath,amssymb,amsthm}
\usepackage{geometry}
\usepackage[colorlinks=true,linkcolor=blue,citecolor=blue,urlcolor=blue]{hyperref}
\usepackage{booktabs}
\usepackage{array}

\geometry{margin=1in}

\newtheorem{theorem}{Theorem}[section]
\newtheorem{proposition}[theorem]{Proposition}
\newtheorem{corollary}[theorem]{Corollary}
\newtheorem{lemma}[theorem]{Lemma}
\theoremstyle{definition}
\newtheorem{definition}[theorem]{Definition}
\theoremstyle{remark}
\newtheorem{remark}[theorem]{Remark}

\title{The Charged-Lepton Spectrum\\[6pt]
\large $S_3$ structure, Koide equipartition, and the generation phase}

\author{Kai Cobbs}
\date{Standalone preprint v0.3\quad May 2026}

\begin{document}
\maketitle

\begin{abstract}
The companion paper \emph{Fisher-Frame Defects and Spinorial Matter} derives
the electron mass $m_e=2m_P\exp(-32\pi^3/19)$ as a single charged spinorial
defect, but does not fix the muon and tau. This paper addresses the
charged-lepton mass \emph{ratios}. The naive route --- successive Dirichlet
eigenvalues of one defect --- fails by two orders of magnitude. Instead the
three generations are identified as the three components of an $S_3$ amplitude
multiplet supported on the $k=3$ (three-prime) sector of the Arithmetic
Vacuum, whose abstract automorphism group is $S_3$. With mass read as squared
amplitude ($m\sim|\psi|^2$, the variable in which the $k=3$ Fisher--Rao metric
linearizes to the round octant of the $2$-sphere), the natural carrier of the
$S_3$ representation is
$\sqrt m$. Two results follow. First, the empirical Koide relation
$Q=(m_e+m_\mu+m_\tau)/(\sqrt{m_e}+\sqrt{m_\mu}+\sqrt{m_\tau})^2 = 0.666661$,
equal to $2/3$ to $1.5\times10^{-5}$, is shown to be \emph{exactly} the
statement that the amplitude sits at the canonical $S_3$ midpoint shell of the
Fisher octant: the colatitude equidistant from the two extremal $S_3$ orbits
(vertices, maximal hierarchy; edge midpoints, maximal pairing), whose
colatitudes are proven to sum to exactly $90^\circ$, forcing the
$45^\circ$/equipartition shell with no free parameter. We further prove an
equivalence: Koide $Q=2/3$ is the vacuum if and only if the $k=3$ defect free
energy is self-dual under the channel duality $I_{\mathbf 1}\leftrightarrow
I_{\mathbf 2}$ (``channels equal'' $\leftrightarrow$ ``channels split''); the
empirical exactness of Koide is then evidence that this duality is exact, and
the dynamical gap is reduced to identifying it (four natural candidate dualities
are excluded). Second, the residual
generation phase is measured to be $\phi=0.222231\pm0.000008$ rad, consistent
with $2/9=(k-1)/k^2$ at $k=3$ to $1.0\sigma$, while the main competing value
$1/(2\sqrt5)$ is excluded at $163\sigma$. Setting $(k,\phi)=(\sqrt2,2/9)$
reproduces all three masses to better than $0.006\%$. We show the phase is a
rigid degeneracy-lifting rotation and prove an obstruction theorem: $\phi=2/9$
\emph{cannot} arise as the vacuum angle of any $S_3$-invariant potential,
because every such invariant depends on the standard-plane angle only through
$\cos(3\psi)$, redirecting the search for its origin to a logarithmic/RG
ratio. The framework predicts, for structural reasons, that leptons are
Koide-exact while quarks are not. Open problems and a near-term falsification
handle (improved $m_\tau$) are stated precisely.
\end{abstract}

\tableofcontents

%==============================================================
\section{The problem and the failure of the naive ladder}
%==============================================================

The companion Fisher-Frame paper produces the electron as a stable charged
$\mathbb{Z}_2$ spinorial defect with mass
\begin{equation}
m_e = 2\,m_P\exp\!\Bigl(-\tfrac{32\pi^3}{19}\Bigr) = 0.5104\ \mathrm{MeV},
\end{equation}
the prefactor $2$ being the first Dirichlet eigenvalue $\lambda_1$ on the
hemisphere. The obvious extension --- assigning the muon and tau to higher
Dirichlet eigenvalues $\lambda\in\{6,12\}$ of the same defect --- gives
$m_\mu/m_e\sim\sqrt{3},\sqrt{6}$, i.e.\ $\sim1.7,2.4$, against the empirical
\begin{equation}
\frac{m_\mu}{m_e}=206.768,\qquad \frac{m_\tau}{m_e}=3477.23.
\end{equation}
The single-defect eigenvalue ladder fails by two orders of magnitude. The
mass law is exponential, so a factor of $207$ is a modest \emph{additive}
shift $\delta_\mu=\ln(m_\mu/m_e)=5.332$ in an exponent of size $\sim52$; the
generations are near-degenerate, split by a sub-leading mechanism, not three
unrelated eigenvalues. This motivates treating them as one multiplet.

%==============================================================
\section{The $S_3$ amplitude multiplet}
%==============================================================

\subsection{Why $\sqrt m$, and why $S_3$}

In the Fisher-frame/symplectic-lift picture, a defect carries an amplitude
$\psi$ and its energy (mass) scales as $m\sim|\psi|^2$ (the same
$\rho=|\psi|^2$ used in the gravity sector's Madelung reduction). The object
transforming linearly under an internal symmetry acting on amplitudes is
therefore $\sqrt m\sim|\psi|$. This is not merely a convention: the $k=3$
sector's state space is the simplex of probability distributions
$p=(p_1,p_2,p_3)$ over the three prime channels, and the Fisher--Rao metric on
the simplex is, in the amplitude coordinates $a_i=\sqrt{p_i}$, exactly the
round metric on the positive octant of the unit $2$-sphere,
$\sum_i a_i^2=1$. The Fisher geometry of the three-prime sector
\emph{linearizes to the sphere precisely in the variable $\sqrt m$}, which is
why $\sqrt m$ --- not $m$ --- is the carrier of the $S_3$ representation. The
three-prime ($k=3$) sector has abstract automorphism group
$\mathrm{Aut}_{\mathrm{abs}}=S_3$ (Physical Sector, Open Problem 11.3), acting
as the symmetry group of the spherical octant. We posit that the three charged
leptons are the three components of the $S_3$ amplitude triplet
\begin{equation}
\boldsymbol a = (\sqrt{m_e},\sqrt{m_\mu},\sqrt{m_\tau})\in\mathbb{R}^3,
\end{equation}
a point of the Fisher octant (up to overall scale), on which $S_3$ acts by
permutation. The democratic axis $(1,1,1)/\sqrt3$ is the octant centre.

\subsection{Representation decomposition}

Under $S_3$, $\mathbb{R}^3=\mathbf 1\oplus\mathbf 2$ (trivial/democratic plus
two-dimensional standard rep), with projections
\begin{equation}
\boldsymbol a_{\mathbf 1}=M(1,1,1),\quad M=\tfrac13\textstyle\sum_n\sqrt{m_n},
\qquad \boldsymbol a_{\mathbf 2}=\boldsymbol a-\boldsymbol a_{\mathbf 1}.
\end{equation}
Numerically (PDG pole masses),
$|\boldsymbol a_{\mathbf 1}|=M\sqrt3=30.6843$ and
$|\boldsymbol a_{\mathbf 2}|=30.6840$, so
$|\boldsymbol a_{\mathbf 2}|/|\boldsymbol a_{\mathbf 1}|=0.999991$.

%==============================================================
\section{Koide as $S_3$ equipartition}
%==============================================================

\subsection{Exact reduction}

Parametrize the amplitude in the general $S_3$-symmetric form
\begin{equation}
\sqrt{m_n}=M\bigl(1+k\cos\phi_n\bigr),\qquad
\phi_n=\theta_0+\tfrac{2\pi n}{3},\quad n=0,1,2,
\label{eq:foot}
\end{equation}
with $M>0$, $k\ge0$. The identities $\sum_n\cos\phi_n=0$ and
$\sum_n\cos^2\phi_n=\tfrac32$ hold for every $\theta_0$, giving
$\sum_n\sqrt{m_n}=3M$ and $\sum_n m_n=3M^2(1+\tfrac{k^2}{2})$.

\begin{proposition}[Koide $=$ equipartition]
\label{prop:koide}
The Koide quotient is
\begin{equation}
Q=\frac{\sum_n m_n}{\bigl(\sum_n\sqrt{m_n}\bigr)^2}=\frac{1+k^2/2}{3},
\end{equation}
so that $Q=\tfrac23$ if and only if $k=\sqrt2$, equivalently
$|\boldsymbol a_{\mathbf 1}|=|\boldsymbol a_{\mathbf 2}|$, equivalently the
amplitude $\boldsymbol a$ makes a $45^\circ$ angle with the democratic axis.
\end{proposition}

\begin{proof}
Direct from the two identities above. For the geometric form,
$\cos^2\theta=|\boldsymbol a_{\mathbf 1}|^2/|\boldsymbol a|^2
=(\sum\sqrt{m_n})^2/(3\sum m_n)=1/(3Q)$, so $Q=2/3\Leftrightarrow
\cos^2\theta=1/2\Leftrightarrow\theta=45^\circ\Leftrightarrow
|\boldsymbol a_{\mathbf 1}|=|\boldsymbol a_{\mathbf 2}|$.
\end{proof}

The empirical value $Q=0.666661$ corresponds to $\theta=44.9997^\circ$. The
entire content of Koide is the single condition $k=\sqrt2$: the
$S_3$-symmetric (singlet) and $S_3$-breaking (standard) parts of the amplitude
have equal norm.

\subsection{What forces $k=\sqrt2$: the canonical midpoint shell}

On the Fisher octant, the equipartition condition has an exact, dynamics-free
characterization in terms of the $S_3$ orbit structure of the boundary. Let
$\Theta$ denote the colatitude of the amplitude measured from the democratic
axis (octant centre), so $I_{\mathbf 1}=\cos^2\Theta$ and
$I_{\mathbf 2}=\sin^2\Theta$, and $k=\sqrt2\Leftrightarrow\Theta=45^\circ$. The
octant boundary carries two distinguished $S_3$ orbits of extremal points:
\begin{itemize}
\item the \emph{vertices} $(1,0,0)$ and permutations --- one prime channel
saturated, maximal generation hierarchy; and
\item the \emph{edge midpoints} $(1,1,0)/\sqrt2$ and permutations --- two
channels equal, one empty.
\end{itemize}

\begin{theorem}[Equipartition is the $S_3$-canonical midpoint shell]
\label{thm:midpoint}
Let $\Theta_v$ and $\Theta_e$ be the colatitudes (from the democratic axis) of
the vertex orbit and the edge-midpoint orbit of the Fisher octant. Then
\begin{equation}
\cos\Theta_v=\tfrac{1}{\sqrt3},\quad \cos\Theta_e=\sqrt{\tfrac23},\qquad
\Theta_v+\Theta_e=\frac{\pi}{2}\ \text{exactly},
\end{equation}
so their midpoint is $\tfrac12(\Theta_v+\Theta_e)=45^\circ$. Hence the
equipartition shell $\Theta=45^\circ$ (Koide $Q=2/3$) is the unique colatitude
equidistant from the two extremal $S_3$ orbits --- the canonical midpoint
shell between maximal hierarchy (vertices) and maximal pairing (edges).
\end{theorem}

\begin{proof}
$\cos\Theta_v=(1,1,1)\!\cdot\!(1,0,0)/\sqrt3=1/\sqrt3$, giving
$\sin\Theta_v=\sqrt{2/3}$. $\cos\Theta_e=(1,1,1)\!\cdot\!(1,1,0)/(\sqrt3\sqrt2)
=\sqrt{2/3}$. Thus $\cos\Theta_e=\sin\Theta_v$, i.e.\
$\Theta_e=\tfrac{\pi}{2}-\Theta_v$, so $\Theta_v+\Theta_e=\tfrac{\pi}{2}$ and
the midpoint is $45^\circ$.
\end{proof}

This upgrades the status of Koide from a tuned coincidence to a structural
characterization: $45^\circ$ is forced by the octant geometry alone, with no
free parameter, as the $S_3$-symmetric halfway shell. Equivalently, defining
the involution $C:\Theta\mapsto\tfrac{\pi}{2}-\Theta$ (acting trivially on the
azimuth), $C$ exchanges the vertex shell with the edge shell, and the
equipartition locus is exactly the fixed shell of $C$. This is the precise
form of the singlet$\leftrightarrow$standard self-duality
($I_{\mathbf 1}\leftrightarrow I_{\mathbf 2}$):

\begin{proposition}[Self-duality form]
\label{prop:selfdual}
The colatitude reflection $C:\Theta\mapsto\tfrac{\pi}{2}-\Theta$ exchanges
$I_{\mathbf 1}\leftrightarrow I_{\mathbf 2}$ and swaps the two extremal $S_3$
orbits; its fixed shell is $\Theta=45^\circ$, i.e.\ $k=\sqrt2$, i.e.\
Koide $Q=2/3$.
\end{proposition}

This is the same type of selection principle as the programme's other balance
points: the self-dual CP angle $\Theta=\pi$; the equal phase-split
$\rho_0=\tfrac12$ between the $\Theta$-term and the spin-$\tfrac12$ Berry
phase; and the spectral balance $g_0^2=J_{\min}^2/S_{\mathrm{esc}}$ that fixes
the gauge coupling.

\subsection{Dynamical selection of the shell}
\label{sec:dynamical}

Theorem~\ref{thm:midpoint} characterizes the Koide shell kinematically. The
remaining question is whether the $k=3$ defect dynamics \emph{selects} it as
the ground state. Here the round geometry of the Fisher--Rao octant plays a
decisive and clarifying role.

\begin{lemma}[Quadratic flatness]
\label{lem:flat}
At quadratic order the colatitude $\Theta$ is a flat direction. The only
quadratic $S_3$-invariant of a unit amplitude is $|\boldsymbol a|^2=1$
(constant on the sphere), so no quadratic potential distinguishes colatitudes;
any selection of $\Theta$ must come from the cubic ($\mathbb{Z}_3$-clock)
invariant or higher.
\end{lemma}

With the azimuth locked by the $\mathbb{Z}_3$ clock (Theorem~\ref{thm:obstruction}
below), an $S_3$-invariant free energy reduces to a function $F(I_{\mathbf 1})$
of the singlet content $I_{\mathbf 1}=\cos^2\Theta\in[\tfrac13,1]$ alone
($I_{\mathbf 1}=1$ at the octant centre, $\tfrac13$ at a vertex, $\tfrac12$ at
the Koide shell). This yields a clean equivalence.

\begin{theorem}[Koide $\iff$ $I_{\mathbf 1}\!\leftrightarrow\!I_{\mathbf 2}$ self-duality of the free energy]
\label{thm:equiv}
An $S_3$-invariant defect free energy $F(I_{\mathbf 1})$ has a stationary point
at the Koide shell $I_{\mathbf 1}=I_{\mathbf 2}=\tfrac12$ if and only if
$F'(\tfrac12)=0$. In particular, if $F$ is invariant under the duality
$C:I_{\mathbf 1}\leftrightarrow I_{\mathbf 2}$ (equivalently $F(I_{\mathbf 1})
=F(1-I_{\mathbf 1})$), then $F'(\tfrac12)=0$ automatically and the shell is
stationary. Conversely, a generic $F$ without this symmetry is stationary
elsewhere (at the centre or a vertex). Hence
\begin{equation}
\text{Koide } Q=\tfrac23 \text{ is the vacuum}
\iff
\text{the } k=3 \text{ free energy is } C\text{-self-dual.}
\end{equation}
\end{theorem}

\begin{proof}
Stationarity at $\tfrac12$ is $F'(\tfrac12)=0$ by definition. If
$F(I_{\mathbf 1})=F(1-I_{\mathbf 1})$ then $F'(I_{\mathbf 1})=-F'(1-I_{\mathbf 1})$,
so $F'(\tfrac12)=-F'(\tfrac12)=0$. The converse direction (genericity) follows
because $F'(\tfrac12)=0$ is one real condition, satisfied automatically on the
$C$-symmetric subspace and only accidentally off it.
\end{proof}

This inverts the logic productively. The empirical exactness of Koide
($Q=2/3$ to $1.5\times10^{-5}$) is, by Theorem~\ref{thm:equiv}, direct
evidence that the $k=3$ defect free energy possesses an \emph{exact}
$C$-duality $I_{\mathbf 1}\leftrightarrow I_{\mathbf 2}$ --- a particle/hole or
occupation duality on the three prime channels, exchanging ``channels equal''
with ``channels maximally split.'' We have therefore reduced the dynamical
question to a single sharp statement: \emph{identify the exact
$C$-duality of the $k=3$ Fisher action.}

\medskip\noindent\textbf{Candidate dualities tested and excluded.} The
following natural involutions all fail to be the required $C$-duality (none has
the $45^\circ$ shell as its fixed/self-dual locus): the global spinorial
$\mathbb{Z}_2$ monodromy (acts trivially on amplitude magnitudes); the
$S_3$-internal transposition (fixes a $2$-plane, not the shell); the
probabilistic complement $p_i\mapsto(1-p_i)/2$ (fixes only the centre); and the
$\mathbb{Z}_3$ discrete Fourier transform (which elegantly swaps centre
$\leftrightarrow$ vertex --- the right \emph{qualitative} behaviour --- but
whose self-dual colatitude is $27.4^\circ$, not $45^\circ$). The variational
restatement is equally clean: several natural functionals ($I_{\mathbf 1}
I_{\mathbf 2}$, the binary entropy $-I_{\mathbf 1}\ln I_{\mathbf 1}
-I_{\mathbf 2}\ln I_{\mathbf 2}$, and $I_{\mathbf 1}^2+I_{\mathbf 2}^2$) are
all extremized at exactly $45^\circ$, but each is manifestly $C$-symmetric, so
each \emph{presupposes} rather than derives the duality.

\medskip\noindent\textbf{Status.} The gap is now sharp and well-posed, not
vague: Theorem~\ref{thm:equiv} reduces ``derive Koide'' to ``exhibit the exact
$C$-duality of the $k=3$ action,'' and the data guarantees such a duality
exists and is exact. Identifying it is the principal open problem
(Section~\ref{sec:open}).

%==============================================================
\section{The generation phase}
\label{sec:phase}
%==============================================================

\subsection{Empirical determination}

After $k=\sqrt2$ is imposed, the spectrum depends on the phase $\theta_0$ and
the scale $M$. Writing $\phi\equiv\theta_0-\tfrac{2\pi}{3}$ and propagating
PDG mass uncertainties by Monte Carlo (dominant source: $m_\tau$, uncertainty
$\sim0.12$ MeV), we obtain
\begin{equation}
\phi=0.222231\pm0.000008\ \mathrm{rad}.
\end{equation}

\begin{center}
\begin{tabular}{lll}\toprule
Candidate & Value (rad) & Deviation\\\midrule
$2/9$ & $0.222222$ & $\mathbf{1.0\,\sigma}$\\
$1/(2\sqrt5)$ (Brannen) & $0.223607$ & $163\,\sigma$ (excluded)\\
$\pi/14$ & $0.224399$ & $271\,\sigma$ (excluded)\\
$7\pi/99$ & $0.222133$ & $12\,\sigma$ (excluded)\\\bottomrule
\end{tabular}
\end{center}

The value $\phi=2/9$ fits at $1.0\sigma$; the main competing closed form is
excluded at $163\sigma$; and no rational multiple of $\pi$ fits within
$2\sigma$. Imposing $(k,\phi)=(\sqrt2,2/9)$ reproduces all three masses to
better than $0.006\%$:

\begin{center}
\begin{tabular}{llll}\toprule
 & predicted (MeV) & actual (MeV) & rel.\ error\\\midrule
$m_e$ & $0.51100$ & $0.51100$ & $<10^{-4}$\\
$m_\mu$ & $105.653$ & $105.658$ & $-0.005\%$\\
$m_\tau$ & $1776.88$ & $1776.86$ & $+0.001\%$\\\bottomrule
\end{tabular}
\end{center}

Structurally $2/9=(k-1)/k^2=2/k^2$ at $k=3$ (these coincide only at $k=3$),
built from the sector index and the standard-rep dimension.

\subsection{What the phase is}

\begin{proposition}[Degeneracy lifting]
At $\phi=0$ the configuration \eqref{eq:foot} with $k=\sqrt2$ has two
degenerate generations, with amplitudes $M(1+\sqrt2)$ and a degenerate pair
$M(1-\tfrac{\sqrt2}{2})$. The physical configuration is the rigid rotation of
the whole triplet by $\phi$; the rotation lifts the two-fold degeneracy.
\end{proposition}

Thus $\phi$ is the degeneracy-splitting parameter of the generation sector,
not a free orientation.

\begin{theorem}[$S_3$ angular obstruction]
\label{thm:obstruction}
The phase $\phi=2/9$ rad cannot arise as the vacuum angle of any
$S_3$-invariant potential. Every $S_3$ invariant depends on the standard-plane
angle $\psi$ only through $\cos(3\psi)$ (the $\mathbb{Z}_3$ clock term
$x^3-3xy^2$), so the stationary points of any $S_3$-invariant potential occur
at $\psi$ equal to integer multiples of $60^\circ$. The observed phase lies
$12.73^\circ=2/9$ rad off the $60^\circ$-lattice.
\end{theorem}

\begin{proof}
The only angular dependence available to an $S_3$-invariant function on the
standard plane is through the real cubic $\mathrm{Re}((x+iy)^3)=
\rho^3\cos(3\psi)$ and its powers; $d/d\psi\,[\cos(3\psi)]=-3\sin(3\psi)$
vanishes precisely at $3\psi\in\pi\mathbb{Z}$, i.e.\ $\psi\in
60^\circ\mathbb{Z}$. The observed standard-plane angle is not on this lattice.
\end{proof}

Theorem~\ref{thm:obstruction} eliminates the most natural mechanism and, taken
with the fact that $\phi$ is a rational number of \emph{radians} (not of
$\pi$, as a geometric angle would be), points the search for the origin of
$2/9$ toward a non-angular logarithmic/RG ratio in the $k=3$ exponent
structure.

\subsection{The democratic scale}

The overall scale is $M^2=313.84$ MeV $=2m_P\exp(-E_M)$ with $E_M=45.80$, a
second dimensional-transmutation scale shifted from the electron exponent
$E_e=52.22$ by $6.42$ (effective one-loop coefficient $b_0^{(M)}=7.22$ versus
the electron's $19/3$). It is not derived here; it requires the explicit
$k=3$-sector beta function and BPS condition, the same calculation type as
Fisher-Frame \S13.

%==============================================================
\section{The quark sector: a structural prediction}
%==============================================================

If equipartition were universal, quarks would also sit at $k=\sqrt2$. They do
not: with running masses the effective values are $k^2_{\mathrm{up}}\approx
3.09$ and $k^2_{\mathrm{down}}\approx2.39$, neither equal to $2$. This is the
known failure of naive Koide for quarks, and here it is consistent with the
framework rather than a problem for it: leptons are the clean $k=3$ spinorial
sector and sit at equipartition, while quarks additionally carry $SU(3)_c$ ---
itself the $k=3$-internal gauge structure (Physical Sector, Open Problem
11.3) --- which distorts the amplitude away from equipartition. The framework
therefore predicts, for a structural reason, that leptons are Koide-exact and
quarks are not, and it offers the quark deviations
$k^2_{\mathrm{up}},k^2_{\mathrm{down}}$ as a future quantitative target
computable from the colour coupling.

%==============================================================
\section{Summary, status, and open problems}
\label{sec:open}
%==============================================================

The charged-lepton spectrum, three free parameters in the Standard Model, is
reduced to:
\begin{enumerate}
\item \textbf{Equipartition $k=\sqrt2$} (Koide $Q=2/3$): exact and
parameter-free, shown to be the canonical $S_3$ midpoint shell of the Fisher
octant --- the colatitude equidistant from the vertex and edge orbits, whose
colatitudes sum to exactly $90^\circ$ (Theorem~\ref{thm:midpoint},
Propositions~\ref{prop:koide}, \ref{prop:selfdual});
\item \textbf{Generation phase $\phi=2/9$}: empirically $1.0\sigma$, the
competing value excluded at $163\sigma$, structurally $(k-1)/k^2$ at $k=3$,
proven not to be an $S_3$ vacuum angle (Theorem~\ref{thm:obstruction});
\item \textbf{Democratic scale $M$}: one undetermined dimensional-transmutation
scale.
\end{enumerate}
This explains the Koide coincidence as representation balance --- something the
Standard Model does not do --- and explains quark non-Koide as $SU(3)_c$
distortion.

\medskip\noindent\textbf{This is a structural characterization plus a sharp
prediction, not yet a closed derivation.} What remains, in priority order:
\begin{enumerate}
\item \textbf{[Critical]} Exhibit the exact $I_{\mathbf 1}\leftrightarrow
I_{\mathbf 2}$ duality $C$ of the $k=3$ Fisher action.
Theorem~\ref{thm:equiv} reduces ``derive Koide'' to exactly this, and the
empirical exactness of $Q=2/3$ guarantees the duality exists and is exact; the
task is to identify it. Four natural candidates have been excluded (global
spinorial $\mathbb{Z}_2$, $S_3$-internal transposition, probabilistic
complement, and the $\mathbb{Z}_3$ DFT --- the last with the correct
centre$\leftrightarrow$vertex behaviour but self-dual colatitude $27.4^\circ$
rather than $45^\circ$), so the duality is a genuine, as-yet-unidentified
structure of the channel dynamics, not one of the obvious involutions.
\item Find $\phi=2/9$ as a logarithmic/RG ratio (not an angle) in the $k=3$
exponent structure; Theorem~\ref{thm:obstruction} shows this is the right
place to look. (Tested and excluded: Shannon channel entropy $0.2186$, and
RG running between pole and high scales, whose shifts exceed the $9\times
10^{-6}$ gap to $2/9$.)
\item Compute the $k=3$ dimensional-transmutation scale $M$ by the
Fisher-Frame \S13 method applied to three primes.
\item Compute the $SU(3)_c$ distortion of quark equipartition and predict
$k^2_{\mathrm{up}},k^2_{\mathrm{down}}$.
\end{enumerate}

\medskip\noindent\textbf{Falsification handle.} Unlike postdictions of known
constants, $\phi=2/9$ is a prediction at the edge of current precision: as
$m_\tau$ is measured more precisely (Belle II, a tau-charm factory), $\phi$
must converge to $0.2222$, not to $0.2236$. A few-fold improvement in the tau
mass will confirm or exclude the result.

\medskip\noindent\textbf{Methodological note.} No quantity was tuned: $k=\sqrt2$
and $\phi=2/9$ are read off the data and tested against structural candidates.
The $1.0\sigma$ versus $163\sigma$ separation is a feature of the data, not a
choice, and the $60^\circ$ obstruction is a genuine negative result that
constrains the mechanism.

%==============================================================
\appendix
\section{Numerical verification}
%==============================================================

All values use PDG central pole masses $m_e=0.51099895$,
$m_\mu=105.6583755$, $m_\tau=1776.86$ MeV.
\begin{itemize}
\item Koide quotient $Q=0.666661$; democratic-axis angle $44.9997^\circ$; vertex colatitude $54.7356^\circ$, edge colatitude $35.2644^\circ$, sum $=90^\circ$ exactly (midpoint $45^\circ$);
$|\boldsymbol a_{\mathbf 2}|/|\boldsymbol a_{\mathbf 1}|=0.999991$.
\item $Q=(1+k^2/2)/3$ verified symbolically.
\item Generic $S_3$ quartic potential: nine couplings; equipartition only on
the self-dual locus.
\item Phase $\phi=0.222231\pm0.000008$ rad ($2\times10^4$ Monte-Carlo samples,
PDG errors, $m_\tau$-dominated).
\item Obstruction: every $S_3$ invariant $\propto\cos(3\psi)$; stationary at
$\psi\in60^\circ\mathbb{Z}$; observed angle off-lattice by $2/9$ rad.
\item Mass reconstruction from $(k=\sqrt2,\phi=2/9,M^2=313.84\,\mathrm{MeV})$
accurate to $<0.006\%$ on all three masses.
\item Flat-direction check: only quadratic $S_3$ invariant is $|a|^2$; colatitude selection requires cubic or higher (Lemma~\ref{lem:flat}).
\item Functionals extremized at $45^\circ$: $I_1 I_2$, binary entropy, $I_1^2+I_2^2$ (all $C$-symmetric). Cubic $e_3=a_1a_2a_3$ and maximal-spread both extremize at the vertex $54.7356^\circ$, not $45^\circ$.
\item $\mathbb{Z}_3$ DFT maps democratic centre to vertex $(1,0,0)$ exactly; its self-dual colatitude is $27.37^\circ\neq45^\circ$.
\item Quark effective values $k^2_{\mathrm{up}}\approx3.09$,
$k^2_{\mathrm{down}}\approx2.39$ (running masses, scheme-dependent).
\end{itemize}

\end{document}
